% ============================================================================
% SECTION 7: NUMERICAL COMPUTATIONS
% ============================================================================
\section{Reproducible Security Diagnostics}\label{sec:numerical}

The computations in this section connect the closed-form geometry to
attack-relevant side-channel diagnostics.  First, they verify the exact
boundary covariance and active-eigenspace formulas on standard lightweight
S-boxes.  Second, they separate quantities that are structural invariants from
quantities that influence CPA discrimination between key hypotheses.  Third,
they validate the pairwise-gap interpretation through a reproducible Monte
Carlo CPA experiment.

\subsection{Target S-boxes}\label{subsec:sboxes}

We evaluate the 4-bit bijective S-boxes used in PRESENT, GIFT, and SKINNY
\cite{bogdanov2007present,banik2017gift,beierle2016skinny}:
\begin{align*}
S_{\mathrm{PRESENT}}&=[C,5,6,B,9,0,A,D,3,E,F,8,4,7,1,2],\\
S_{\mathrm{GIFT}}&=[1,A,4,C,6,F,3,9,2,D,B,7,5,0,8,E],\\
S_{\mathrm{SKINNY}}&=[C,6,9,0,1,A,2,B,3,8,5,D,4,E,7,F].
\end{align*}
These examples are useful because they are widely used in lightweight
cryptography and because exhaustive enumeration over all inputs, masks, and
key hypotheses is possible.

\subsection{Walsh spectra and boundary covariance}\label{subsec:wht-spectra}

For each S-box we compute the full WHT table
\[
  \widehat S(\alpha,\beta)=\sum_{a\in\B{4}}
  (-1)^{\alpha\cdot a\oplus\beta\cdot S(a)}.
\]
The three tested S-boxes have the same absolute Walsh-coefficient distribution
over nonzero masks:
\[
  \#|\widehat S|=0:123,\qquad
  \#|\widehat S|=4:96,\qquad
  \#|\widehat S|=8:36.
\]
Therefore diagnostics depending only on the Walsh amplitudes coincide for
these three examples.  This shared spectral profile provides a useful baseline
for reading the later CPA-specific quantities.

\begin{figure}[t]
\centering
\includegraphics[width=\textwidth]{fig1_wht_spectra.pdf}
\caption{Walsh spectra. The plotted Walsh tables show the same amplitude
profile for PRESENT, GIFT, and SKINNY, which explains why purely
amplitude-based diagnostics coincide for these examples.}
\label{fig:wht-spectra}
\end{figure}

\begin{figure}[t]
\centering
\includegraphics[width=0.74\textwidth]{fig2_wht_histogram.pdf}
\caption{Walsh amplitude distribution. The histogram records the shared counts
of absolute Walsh coefficients over nonzero masks for the three tested S-boxes.}
\label{fig:wht-histogram}
\end{figure}

The boundary covariance matrix is built from Proposition~\ref{thm:fim}.  For
$s=4$ it is a $255\times255$ positive semidefinite matrix.  Consistent with
Theorem~\ref{thm:active-rank}, all three matrices have rank $15$ and active
eigenvalues equal to $16$ up to numerical precision.

\begin{table}[t]
\centering
\caption{Boundary covariance structure for the tested 4-bit bijections.}
\label{tab:cov-rank}
\begin{tabular}{lccc}
\toprule
S-box & Matrix size & Rank & Active eigenvalues \\
\midrule
PRESENT & $255\times255$ & $15$ & $16$ repeated $15$ times \\
GIFT    & $255\times255$ & $15$ & $16$ repeated $15$ times \\
SKINNY  & $255\times255$ & $15$ & $16$ repeated $15$ times \\
\bottomrule
\end{tabular}
\end{table}

\subsection{Local spectral separation}\label{subsec:fr-distances}

Using~\eqref{eq:fr-local}, we compute local diagonal Walsh/Fisher separation
proxies between key hypotheses.  This local covariance-weighted coordinate
displacement is evaluated in the regularized Walsh/Fisher coordinates
introduced earlier.  The key-translation rule implies that the value depends
only on the key difference $k\oplus k'$.

\begin{figure}[t]
\centering
\includegraphics[width=0.8\textwidth]{fig3_fr_distance_matrix.pdf}
\caption{Local key-separation proxy. The matrix visualises the
covariance-weighted Walsh displacement between key hypotheses in the local
Walsh/Fisher coordinates.}
\label{fig:local-distance}
\end{figure}

\begin{figure}[t]
\centering
\includegraphics[width=0.74\textwidth]{fig5_fr_vs_hw.pdf}
\caption{Local proxy versus key-difference Hamming weight.  The plot aggregates
the same covariance-weighted Walsh displacement by the Hamming weight of
$k\oplus k'$.}
\label{fig:local-distance-hw}
\end{figure}

\begin{table}[t]
\centering
\caption{Interpretation of the key-separation quantities.}
\label{tab:distance-status}
\begin{tabular}{ll}
\toprule
Quantity & Status \\
\midrule
$\KL(P_k\|P_{k'})$ for $k\neq k'$ & Infinite, because supports are disjoint \\
$d_F(\delta_k,\delta_{k'})$ on the belief simplex & Exactly $\pi$ \\
$d_{\mathrm{loc}}^2(k,k')$ in~\eqref{eq:fr-local} & Local Walsh/Fisher covariance proxy \\
\bottomrule
\end{tabular}
\end{table}

\subsection{CPA leakage gaps}\label{subsec:pairwise-gap-numerical}

For bijective 4-bit S-boxes and uniform inputs, the scalar Hamming-weight
variance in Proposition~\ref{thm:per-trace} is always
\[
  I_S(\HW,\sigma)=\frac{1}{\sigma^2}.
\]
This scalar is a baseline leakage descriptor shared by PRESENT, GIFT, and
SKINNY.  Pairwise gaps
\[
  \Delta(\delta)=\frac1{16}\sum_{a\in\B{4}}
  \bigl(\HW(S(a))-\HW(S(a\oplus\delta))\bigr)^2,
  \qquad \delta\neq0,
\]
are more informative for CPA because the expected log-Bayes gain against a
wrong key with difference $\delta$ is $\Delta(\delta)/(2\sigma^2)$.

\begin{table}[t]
\centering
\caption{Pairwise Hamming-weight leakage gaps over nonzero key differences.}
\label{tab:pairwise-gaps}
\begin{tabular}{lccc}
\toprule
S-box & $\min_{\delta\neq0}\Delta(\delta)$ &
$\max_{\delta\neq0}\Delta(\delta)$ &
$\frac1{15}\sum_{\delta\neq0}\Delta(\delta)$ \\
\midrule
PRESENT & $1$ & $7/2$ & $32/15$ \\
GIFT    & $5/4$ & $7/2$ & $32/15$ \\
SKINNY  & $1$ & $3$   & $32/15$ \\
\bottomrule
\end{tabular}
\end{table}

Table~\ref{tab:pairwise-gaps} shows how the scalar variance descriptor and the
pairwise structure play complementary roles: all three S-boxes have the same
average pairwise gap, while the worst and best key-difference cases vary.
These pairwise gaps are the
quantities that enter expected posterior-odds growth in the Gaussian leakage
model.

\begin{table}[t]
\centering
\caption{Operational CPA interpretation of the pairwise gaps.  The listed
deltas are nonzero key differences $\delta=k\oplus k^*$ in hexadecimal.}
\label{tab:gap-distribution}
\scriptsize
\setlength{\tabcolsep}{2pt}
\begin{tabular}{@{}p{0.13\textwidth}p{0.17\textwidth}p{0.17\textwidth}p{0.43\textwidth}@{}}
\toprule
S-box & Hardest gaps $\min\Delta$ & Easiest gaps $\max\Delta$ &
Gap multiset over $\delta\neq0$ \\
\midrule
PRESENT & $1$ at $\{A,E\}$ & $7/2$ at $\{6\}$ &
$1^2$; $(5/4)^1$; $(3/2)^2$; $(7/4)^3$; $(5/2)^2$;
$(11/4)^1$; $3^2$; $(13/4)^1$; $(7/2)^1$ \\
GIFT & $5/4$ at $\{C,E,F\}$ & $7/2$ at $\{7\}$ &
$(5/4)^3$; $(3/2)^2$; $2^2$; $(9/4)^3$; $(5/2)^2$;
$3^2$; $(7/2)^1$ \\
SKINNY & $1$ at $\{5\}$ & $3$ at $\{B\}$ &
$1^1$; $(5/4)^1$; $(3/2)^1$; $(7/4)^1$; $2^3$;
$(9/4)^2$; $(5/2)^3$; $(11/4)^2$; $3^1$ \\
\bottomrule
\end{tabular}
\end{table}

Table~\ref{tab:gap-distribution} gives the CPA-oriented reading of the
pairwise theory.  PRESENT and SKINNY have a smaller hardest-key gap than GIFT
under this leakage model, so the model-based estimate
\eqref{eq:gap-sample-estimate} assigns a larger trace budget to the hardest
posterior-odds margin for those S-boxes.  The comparison is local to the
specified S-box computation and leakage model, which is the level at which CPA
key hypotheses are evaluated.

\subsection{Monte Carlo CPA validation}\label{subsec:cpa-simulation}

To connect the gap analysis to an attack observable, we simulate the Bayesian
CPA model of Section~\ref{sec:cpa}.  For each S-box, we use true key
$k^*=0$, uniformly sampled plaintext nibbles, leakage
$W=\HW(S(a\oplus k^*))+\varepsilon$, Gaussian noise standard deviation
$\sigma=2$, and $3000$ Monte Carlo trials with seed $20260622$.  For every
trace prefix, all 16 key log-likelihoods are accumulated and the true-key rank
is recorded.  The script \texttt{simulate\_cpa.py} regenerates both
Fig.~\ref{fig:cpa-simulation} and the CSV summary
\texttt{cpa\_simulation\_results.csv}.

\begin{figure}[t]
\centering
\includegraphics[width=\textwidth]{fig6_cpa_simulation.pdf}
\caption{Simulated CPA under Hamming-weight Gaussian leakage.  The left panel
shows empirical success probability and the right panel shows mean true-key
rank over $3000$ trials.  The curves translate the pairwise-gap analysis into
trace-dependent attack observables.}
\label{fig:cpa-simulation}
\end{figure}

The simulation confirms the qualitative role of the pairwise gaps.  At
$64$ traces, the empirical success probabilities are approximately $0.923$ for
PRESENT, $0.950$ for GIFT, and $0.947$ for SKINNY; at $128$ traces, all three
success probabilities are between $0.993$ and $0.996$.  Mean true-key ranks
decrease toward one for all three S-boxes.  The finite-trace curves summarize
the combined effect of all wrong-key competitors, while the minimum gap in
Table~\ref{tab:pairwise-gaps} controls the hardest pairwise posterior-odds
growth.

\begin{figure}[t]
\centering
\includegraphics[width=0.74\textwidth]{fig4_per_trace_info.pdf}
\caption{Per-trace leakage baseline. The scalar Hamming-weight variance is the
same for all bijective 4-bit S-boxes under uniform inputs, so CPA discrimination
is completed by pairwise leakage gaps that determine key-hypothesis
discrimination.}
\label{fig:per-trace-info}
\end{figure}

\subsection{Active-eigenspace curvature}\label{subsec:curv-numerical}

The curvature diagnostic of Section~\ref{sec:curvature} is evaluated in the
15-dimensional active eigenspace of the boundary covariance matrix.  The result
is constant for all coordinate planes in the tested cases.

\begin{table}[t]
\centering
\caption{Projected active-eigenspace curvature diagnostic for the tested 4-bit S-boxes.}
\label{tab:curvatures}
\begin{tabular}{lcccc}
\toprule
S-box & Active dimension & Planes & Min $K$ & Max $K$ \\
\midrule
PRESENT & $15$ & $105$ & $-0.25$ & $-0.25$ \\
GIFT    & $15$ & $105$ & $-0.25$ & $-0.25$ \\
SKINNY  & $15$ & $105$ & $-0.25$ & $-0.25$ \\
\bottomrule
\end{tabular}
\end{table}

\subsection{Reproducibility checklist}\label{subsec:repro-checklist}

The computational claims in this section are reproducible from the additional
files.  The scripts \texttt{gen\_figures.py}, \texttt{simulate\_cpa.py},
\texttt{compute\_curvatures.py}, and \texttt{verify\_curvature.py} regenerate
the figures, CPA simulation results, active covariance diagnostics, and
curvature checks.  The
rank/eigenvalue statement is proved independently in
Theorem~\ref{thm:active-rank}; the scripts provide finite checks for the
reported 4-bit examples.
